\documentclass[10pt,a4paper]{amsart}
\usepackage{amsmath,amssymb,amsthm,graphicx}
% Retained minimal preamble for the Stacks Project chapter more-algebra.tex (More on Algebra).
% The chapter file pinned-original-source.tex carries only \input{preamble} on its first line, so no
% preamble text lives inside the chapter itself.  The Stacks Project uses one single shared preamble for
% every chapter; the snapshot stages it identically as sources/stacks-chow/preamble.tex and
% sources/stacks-etale-cohomology/preamble.tex.  The block below is copied mechanically, line for line,
% from sources/stacks-chow/preamble.tex lines 14-18 (the Xy-pic setup) and 154-217 (the native theorem
% declarations together with the mathematical macros).  Nothing was added, removed or redefined; the
% wrapper class and the amsmath/amssymb/amsthm/graphicx base packages are supplied by the assembler.
% Chapter: more-algebra.tex -- More on Algebra.
% For commutative diagrams we use Xy-pic
\usepackage[all]{xy}

% We use 2cell for 2-commutative diagrams.
\xyoption{2cell}

% Theorem environments.
%
\theoremstyle{plain}
\newtheorem{theorem}[subsection]{Theorem}
\newtheorem{proposition}[subsection]{Proposition}
\newtheorem{lemma}[subsection]{Lemma}

\theoremstyle{definition}
\newtheorem{definition}[subsection]{Definition}
\newtheorem{example}[subsection]{Example}
\newtheorem{exercise}[subsection]{Exercise}
\newtheorem{situation}[subsection]{Situation}

\theoremstyle{remark}
\newtheorem{remark}[subsection]{Remark}
\newtheorem{remarks}[subsection]{Remarks}

\numberwithin{equation}{subsection}

% Macros
%
\def\lim{\mathop{\mathrm{lim}}\nolimits}
\def\colim{\mathop{\mathrm{colim}}\nolimits}
\def\Spec{\mathop{\mathrm{Spec}}}
\def\Hom{\mathop{\mathrm{Hom}}\nolimits}
\def\Ext{\mathop{\mathrm{Ext}}\nolimits}
\def\SheafHom{\mathop{\mathcal{H}\!\mathit{om}}\nolimits}
\def\SheafExt{\mathop{\mathcal{E}\!\mathit{xt}}\nolimits}
\def\Sch{\mathit{Sch}}
\def\Mor{\mathop{\mathrm{Mor}}\nolimits}
\def\Ob{\mathop{\mathrm{Ob}}\nolimits}
\def\Sh{\mathop{\mathit{Sh}}\nolimits}
\def\NL{\mathop{N\!L}\nolimits}
\def\CH{\mathop{\mathrm{CH}}\nolimits}
\def\proetale{{pro\text{-}\acute{e}tale}}
\def\etale{{\acute{e}tale}}
\def\QCoh{\mathit{QCoh}}
\def\Ker{\mathop{\mathrm{Ker}}}
\def\Im{\mathop{\mathrm{Im}}}
\def\Coker{\mathop{\mathrm{Coker}}}
\def\Coim{\mathop{\mathrm{Coim}}}

% Boxtimes
%
\DeclareMathSymbol{\boxtimes}{\mathbin}{AMSa}{"02}

%
% Macros for moduli stacks/spaces
%
\def\QCohstack{\mathcal{QC}\!\mathit{oh}}
\def\Cohstack{\mathcal{C}\!\mathit{oh}}
\def\Spacesstack{\mathcal{S}\!\mathit{paces}}
\def\Quotfunctor{\mathrm{Quot}}
\def\Hilbfunctor{\mathrm{Hilb}}
\def\Curvesstack{\mathcal{C}\!\mathit{urves}}
\def\Polarizedstack{\mathcal{P}\!\mathit{olarized}}
\def\Complexesstack{\mathcal{C}\!\mathit{omplexes}}
% \Pic is the operator that assigns to X its picard group, usage \Pic(X)
% \Picardstack_{X/B} denotes the Picard stack of X over B
% \Picardfunctor_{X/B} denotes the Picard functor of X over B
\def\Pic{\mathop{\mathrm{Pic}}\nolimits}
\def\Picardstack{\mathcal{P}\!\mathit{ic}}
\def\Picardfunctor{\mathrm{Pic}}
\def\Deformationcategory{\mathcal{D}\!\mathit{ef}}
\usepackage{cleveref}
\crefname{originallemma}{original lemma}{original lemmas}
\Crefname{originallemma}{Original lemma}{Original lemmas}
\makeatletter
\newlabel{lemma-last-one-flat}{{0653}{source}{}{sourcepointer-1}{}}
\newlabel{lemma-last-one-flat@cref}{{[originallemma][0][]0653}{[1][0][]source}{}{}{}}
\makeatother
\crefname{originallemma}{original lemma}{original lemmas}
\Crefname{originallemma}{Original lemma}{Original lemmas}
\makeatletter
\newlabel{lemma-RHom-out-of-projective}{{0A66}{source}{}{sourcepointer-2}{}}
\newlabel{lemma-RHom-out-of-projective@cref}{{[originallemma][0][]0A66}{[1][0][]source}{}{}{}}
\makeatother
\crefname{originallemma}{original lemma}{original lemmas}
\Crefname{originallemma}{Original lemma}{Original lemmas}
\makeatletter
\newlabel{lemma-base-change-comparison}{{0661}{source}{}{sourcepointer-3}{}}
\newlabel{lemma-base-change-comparison@cref}{{[originallemma][0][]0661}{[1][0][]source}{}{}{}}
\makeatother
\title{An explicit R-flat complex computing RHom between a pseudo-coherent object and an object of finite tor dimension, and its base change}
\author{The Stacks Project Authors}
\date{stacks-project:more-algebra.tex:local-snapshot-20261012; GNU FDL 1.2 or later}
\begin{document}
\maketitle
\setcounter{subsection}{1}
\noindent\textit{Editorial note.} All delivered mathematical text is the author's own, copied from the pinned Stacks Project chapter file more-algebra.tex; no statement or proof step was rewritten, reordered or removed. The counted statement is the labeled lemma environment at source lines 22039-22062 (canonical tag 0E1U) and the counted proof is the proof environment at source lines 22064-22102, which immediately follows it and is its own proof.\par\medskip
\noindent\textit{Editorial note.} The printed text contains the slips since F is bounded an since P is bounded above at source line 22076, where and is meant. The text is copied exactly as it stands and is not corrected.\par\medskip
\noindent\textit{Editorial note.} The three claims of the statement that are not proved in the span itself are the two choices of representatives and the finite presentation refinement of the last two sentences; the first is the definition of pseudo-coherence, the second is derived from bounded tor dimension together with the truncation and the cited flatness statement, and the refinement is a choice that the preceding arguments support, since a relatively perfect object admits finitely presented representatives.\par\medskip
\noindent\textit{Editorial note.} The chapter file pinned-original-source.tex carries only an input of the shared preamble on its first line, so no preamble text lives inside the chapter. The retained file sources/stacks-more-algebra/retained-preamble.tex is a mechanical, line-for-line copy of the project's single shared preamble as staged in the snapshot (Xy-pic setup plus the native theorem declarations and mathematical macros). Nothing was added, removed or redefined, and the excerpt serves only as preamble, never as a span coordinate.\par\medskip
\noindent\textit{Editorial note.} The source URL is the source-level Stacks Project URL, as required when more than one card is produced from one source; the per-claim tag is carried by claim\_key and by this note. The counted statement and proof cite no bibliography entry, so this card needs no bibliography.\par\medskip
\noindent\textit{Allowed standard prerequisites (editorial).}\begin{itemize}
\item A pseudo-coherent complex is by definition represented by a bounded above complex of finite free modules, and a bounded tor dimension implies boundedness, so a bounded above complex of free modules represents the object and can be truncated. used at proof lines 22065-22072. Produces the two complexes P and F of the statement and makes F a bounded complex of R-flat modules.
\item For a bounded above complex of finite free modules and a bounded complex of flat modules the Hom complex is a bounded below complex whose terms are finite direct sums of the Hom modules in the matching degrees. used at proof lines 22075-22082. Gives the termwise description of E and shows that its terms are R-flat, which is assertion (a) apart from representing RHom.
\item Base change for RHom along R -> R' and the compatibility of Hom with base change for the terms, using flatness of the terms over R. used at proof lines 22087-22101. Identifies E tensor\_R R' with the Hom complex of the base changed complexes, which is assertion (b).
\end{itemize}\medskip
\ifdefined\hypertarget\hypertarget{sourcepointer-1}{}\fi
\noindent\textit{Source reference.} Original lemma 0653: sources/stacks-more-algebra/pinned-original-source.tex, source lines 16883--16883. The reference is used at proof line 22072 to see that the truncation of F has R-flat terms; the termwise computation of the Hom complex that follows is inside the counted span.\par\medskip
\ifdefined\hypertarget\hypertarget{sourcepointer-2}{}\fi
\noindent\textit{Source reference.} Original lemma 0A66: sources/stacks-more-algebra/pinned-original-source.tex, source lines 18708--18708. The reference identifies the complex E with RHom\_A(K, L) at proof lines 22083-22084; the termwise description of E that makes it a bounded below complex of R-flat modules is inside the counted span.\par\medskip
\ifdefined\hypertarget\hypertarget{sourcepointer-3}{}\fi
\noindent\textit{Source reference.} Original lemma 0661: sources/stacks-more-algebra/pinned-original-source.tex, source lines 14955--14955. The reference is used at proof lines 22089-22092 to see that the base changed complexes represent the base changed objects; the subsequent termwise identification of E tensor\_R R' with the Hom complex of those representatives is inside the counted span.\par\medskip
\begin{lemma}
\label{lemma-compute-RHom-relatively-perfect}
Let $R \to A$ be a flat ring map. Let $K, L \in D(A)$ with $K$
pseudo-coherent and $L$ finite tor dimension over $R$. We may choose
\begin{enumerate}
\item a bounded above complex $P^\bullet$
of finite free $A$-modules representing $K$, and
\item a bounded complex of $R$-flat $A$-modules
$F^\bullet$ representing $L$.
\end{enumerate}
Given these choices we have
\begin{enumerate}
\item[(a)] $E^\bullet = \Hom^\bullet(P^\bullet, F^\bullet)$
is a bounded below complex
of $R$-flat $A$-modules representing $R\Hom_A(K, L)$,
\item[(b)] for any ring map $R \to R'$ with $A' = A \otimes_R R'$
the complex $E^\bullet \otimes_R R'$ represents
$R\Hom_{A'}(K \otimes_A^\mathbf{L} A', L \otimes_A^\mathbf{L} A')$.
\end{enumerate}
If in addition $R \to A$ is of finite presentation and $L$ is
$R$-perfect, then we may choose $F^p$ to be finitely presented
$A$-modules and consequently $E^n$ will be finitely presented $A$-modules
as well.
\end{lemma}

\begin{proof}
The existence of $P^\bullet$ is the definition of a pseudo-coherent complex.
We first represent $L$
by a bounded above complex $F^\bullet$ of free $A$-modules
(this is possible because bounded tor dimension in particular
implies bounded). Next, say $L$ viewed as an object of $D(R)$
has tor amplitude in $[a, b]$. Then, after replacing
$F^\bullet$ by $\tau_{\geq a}F^\bullet$, we get a complex
as in (2). This follows from Lemma \ref{lemma-last-one-flat}.

\medskip\noindent
Proof of (a).
Since $F^\bullet$ is bounded an since $P^\bullet$ is bounded above,
we see that $E^n = 0$ for $n \ll 0$ and that
$E^n$ is a finite (!) direct sum
$$
E^n = \bigoplus\nolimits_{p + q = n} \Hom_A(P^{-q}, F^p)
$$
and since $P^{-q}$ is finite free, this is indeed an $R$-flat $A$-module.
The fact that $E^\bullet$ represents $R\Hom_A(K, L)$
follows from Lemma \ref{lemma-RHom-out-of-projective}.

\medskip\noindent
Proof of (b).
Let $R \to R'$ be a ring map and $A' = A \otimes_R R'$.
By Lemma \ref{lemma-base-change-comparison} the object
$L \otimes_A^\mathbf{L} A'$ is represented by
$F^\bullet \otimes_R R'$ viewed as a complex of $A'$-modules
(by flatness of $F^p$ over $R$). Similarly for $P^\bullet \otimes_R R'$.
As above $R\Hom_{A'}(K \otimes_A^\mathbf{L} A', L \otimes_A^\mathbf{L} A')$
is represented by
$$
\Hom^\bullet(P^\bullet \otimes_R R', F^\bullet \otimes_R R') =
E^\bullet \otimes_R R'
$$
The equality holds by looking at the terms of the complex individually
and using that $\Hom_{A'}(P^{-q} \otimes_R R', F^p \otimes_R R') =
\Hom_A(P^{-q}, F^p) \otimes_R R'$.
\end{proof}

\end{document}
